% Transcription — fonds Alexandre Grothendieck, shelfmark 160, batch 2 (pages 21-22).
% Established from the facsimile archives/batches/160/batch-02.pdf (a mirror of
% https://grothendieck.umontpellier.fr/160.pdf), per the skill
% transcribe-grothendieck. The batch file carries no cover sheet and does not
% start at page 1: PDF page 1 = archivists' page 21, PDF page 2 = page 22,
% checked against the pencilled numbers bottom-left ("21" and "22" both read).
%
% Pass: Opus 5 (claude-opus-5), 2026-09-19 — first pass, unchecked against the pages by a human.
%
% What the batch is. Two leaves, the last of the folder, and both continue the
% run of batch 1 without a break — page 20 ends inside the same calculation
% that page 21 takes up, on additive polynomials X^{p^2} + aX^p + bX over a
% field of characteristic p and on the finite subgroup schemes of G_a they cut
% out. Neither leaf carries a sentence of connected prose: page 21 is pure
% calculation, page 22 is a page of schematic notes (inclusions, products,
% two small diagrams) with abbreviations for connective tissue.
%
% Page 22 is written from both ends: the upper third upright, the lower two
% thirds head to foot, the two separated by a long pencil rule. That is the
% author's own layout, not a scanning accident (the archivists' "22" is
% upright), and it fixes the reading order: the inverted block is read from
% the foot of the sheet upwards, and is transcribed after the upright block.
% Below its rule the inverted block is itself in two columns, divided by a
% vertical stroke; the left column is given first, then the right.

\documentclass[11pt,a4paper]{article}
% Le préambule commun à toutes les transcriptions du fonds Grothendieck.
%
% Il tient en une idée : **l'incertitude se compose, elle ne se résout pas.**
% Une transcription qui lisse un mot douteux a détruit la seule chose pour
% laquelle on la faisait — un lecteur doit pouvoir savoir, à chaque endroit,
% ce qui était sur le feuillet et ce qui a été ajouté. Les six macros qui
% suivent sont donc l'appareil critique, et non de la décoration.
%
% Les mêmes commandes sont reconnues par `scripts/render.mjs`, qui produit la
% vue de lecture du site. Une macro ajoutée ici doit l'être aussi là-bas,
% faute de quoi le rendu échouera bruyamment — ce qui est voulu : un rendu
% silencieusement faux est pire qu'un rendu absent.

\ProvidesPackage{grothendieck}[2026/08/08 Transcriptions du fonds Grothendieck]

\usepackage{amsmath,amssymb,amsthm}
\usepackage{mathrsfs}
\usepackage{stmaryrd}
% Les diagrammes commutatifs sont partout dans ce fonds : c'est la langue
% même de la géométrie algébrique de Grothendieck, et une transcription qui
% les rendrait en prose ne transcrirait rien.
\usepackage{tikz-cd}
% Français par défaut — c'est la langue des feuillets — mais l'anglais est
% chargé aussi : la traduction et le résumé sont des documents anglais, et la
% typographie française (espaces avant les deux-points) y serait une faute.
% Ces documents basculent par \selectlanguage{english} en tête de corps.
\usepackage[english,french]{babel}
\usepackage{fontspec}
\usepackage{geometry}
\geometry{a4paper,margin=25mm,marginparwidth=28mm}
\usepackage{marginnote}
\usepackage{xcolor}
% ulem, not soul. `soul`'s \st and \ul are built on a character-by-character
% scan that XeLaTeX's font machinery breaks: the document fails with an
% undefined control sequence rather than degrading. ulem does the same two jobs
% and is Unicode-engine safe. `normalem` stops it hijacking \emph, which the
% transcriptions use for Grothendieck's own emphasis.
\usepackage[normalem]{ulem}
\usepackage{hyperref}
\hypersetup{colorlinks=true,linkcolor=black,urlcolor=black,pdfborder={0 0 0}}

% --- Métadonnées du lot -----------------------------------------------------
% Reprises telles quelles par le script de rendu, qui les affiche en tête de
% la vue de lecture. Elles fixent ce que le fichier prétend couvrir : sans
% elles, une transcription flotte sans point d'attache dans un fonds de seize
% mille feuillets.

\newcommand{\folder}[1]{\gdef\grothFolder{#1}}
\newcommand{\batch}[1]{\gdef\grothBatch{#1}}
\newcommand{\leaves}[2]{\gdef\grothFirst{#1}\gdef\grothLast{#2}}
\newcommand{\foldertitle}[1]{\gdef\grothFoldertitle{#1}}
\gdef\grothFolder{?}\gdef\grothBatch{?}\gdef\grothFirst{?}\gdef\grothLast{?}\gdef\grothFoldertitle{}

% --- Le résumé --------------------------------------------------------------
%
% Il ouvre la lecture modernisée, sous le titre « Résumé », et il est composé
% en petit. Ce n'est pas un ornement : c'est le seul passage du document qui
% ne soit pas des mathématiques, et un lecteur doit pouvoir le reconnaître
% avant de l'avoir lu — puis le sauter s'il connaît déjà le sujet, ou s'y
% arrêter s'il ne le connaît pas. Le filet qui le suit marque la reprise du
% corps.

\newenvironment{resume}
  {\small\setlength{\parskip}{0.45em}}
  {\par\medskip\hrule\medskip}

% --- Mots-clés ---------------------------------------------------------------
%
% En anglais, en clôture du résumé de la lecture modernisée : le vocabulaire
% moderne sous lequel chercher ce que les feuillets construisent. Le manifeste
% du site (`npm run manifest`) les extrait pour étiqueter la cote entière —
% cette ligne est la source unique des tags, il n'y a pas de fichier à côté.
\newcommand{\keywords}[1]{\par\smallskip\noindent
  {\footnotesize\textsc{Keywords} --- \textit{#1}}\par}

% --- L'appareil critique ----------------------------------------------------

% Le numéro de feuillet, en marge. C'est l'ancre qui permet de régler un
% désaccord sur une formule en regardant *un* feuillet plutôt qu'une chemise
% de six cents. Le site s'en sert aussi pour tourner les pages du fac-similé
% quand on fait défiler la transcription.
\newcommand{\leaf}[1]{%
  \marginnote{\footnotesize\color{gray!70}\textsf{#1}}%
  \label{leaf:#1}\ignorespaces}

% Illisible : on ne devine pas. Un mot inventé qui se lit comme les autres est
% le pire résultat possible.
\newcommand{\ill}{\textcolor{red!60!black}{[\,\ldots\,]}}

% Lecture douteuse : le mot est donné, et signalé comme incertain.
\newcommand{\uncertain}[1]{\uline{#1}}

% Ajout de l'éditeur — une lettre manquante, un mot sous-entendu.
\newcommand{\add}[1]{\textcolor{blue!50!black}{[#1]}}

% Biffé par Grothendieck lui-même. Ce qu'il a rayé fait partie du document :
% une rature dit souvent où la pensée a changé de direction.
\newcommand{\struck}[1]{\sout{#1}}

% Note du transcripteur, sur l'état matériel du feuillet.
\newcommand{\note}[1]{\par\smallskip{\small\color{gray!60!black}\textit{[#1]}}\par\smallskip}

% Note marginale de Grothendieck, distincte du corps du texte.
\newcommand{\marginal}[1]{\par\smallskip\noindent\hspace{1em}%
  {\small\color{gray!60!black}#1}\par\smallskip}

% --- Titre ------------------------------------------------------------------

\AtBeginDocument{%
  \begin{center}
    {\large\bfseries Fonds Alexandre Grothendieck --- cote n\textsuperscript{o}~\grothFolder}\\[2pt]
    {\itshape\grothFoldertitle}\\[4pt]
    {\small feuillets \grothFirst--\grothLast\ (lot \grothBatch)}\\[2pt]
    {\footnotesize\color{gray!60!black}Transcription établie d'après le fac-similé numérisé
     par l'Université de Montpellier.\\
     Les lectures incertaines sont soulignées, les illisibles notées [\,\ldots\,],
     les ajouts entre crochets.}
  \end{center}
  \vspace{1em}}

\endinput


\folder{160}
\batch{2}
\pages{21}{22}
\dating{s.d. — le groupe « Dérivateurs » (157-1 à 160) est daté 1990-[1991]}
\watermark{Édition de démonstration}
\foldertitle{Schémas [en] groupes : notes manuscrites (s.d.)}

\begin{document}

\page{21}
\note{feuillet de calcul, sans une phrase ; il reprend sans coupure le calcul
du batch 1 sur les polynômes additifs en caractéristique $p$}

$\alpha$ \qquad $\beta$

\[
  X^{p^{2}} + aX^{p} + bX
\]
\note{$\alpha$ et $\beta$ sont portés au-dessus du polynôme : ce sont les deux
éléments dont part l'élimination qui suit}

\begin{align*}
  \alpha^{p^{2}} + a\alpha^{p} + b\alpha &= 0\\
  \beta^{p^{2}} + a\beta^{p} + b\beta &= 0
\end{align*}

\[
  a\bigl(\alpha^{p}\beta - \beta^{p}\alpha\bigr)
  = -\bigl(\alpha^{p^{2}}\beta - \beta^{p^{2}}\alpha\bigr)
\]

\[
  a = \frac{\alpha^{p^{2}}\beta - \beta^{p^{2}}\alpha}
           {\alpha^{p}\beta - \beta^{p}\alpha}
  \qquad
  b = \frac{\alpha^{p^{2}}\beta^{p} - \beta^{p^{2}}\alpha^{p}}
           {\alpha^{p}\beta^{\struck{p}} - \beta^{p}\alpha^{\struck{p}}}
\]
\note{au dénominateur de $b$ les deux exposants $p$ portés sur $\beta$ puis sur
$\alpha$ sont biffés : le dénominateur devient celui de $a$}

\note{la colonne de droite du feuillet refait le même calcul en se servant de
$\alpha^{p}-\beta^{p} = (\alpha-\beta)^{p}$ ; elle est donnée ici à la suite}

\[
  a\bigl(\alpha^{p} - \beta^{p}\bigr)
  = -\bigl(\alpha^{p^{2}} - \beta^{p^{2}}\bigr)
\]

\begin{align*}
  a &= -(\alpha-\beta)^{p^{2}-p}\\
  b &= -\alpha^{p^{2}} + \alpha^{p}(\alpha-\beta)^{p^{2}-p}
\end{align*}
\note{le premier membre de la seconde ligne est écrit $b$ : la relation tirée
de $\alpha^{p^{2}} + a\alpha^{p} + b\alpha = 0$ est celle de $b\alpha$. Le même
facteur $\alpha$ manque au premier membre des deux lignes suivantes, qui sont
cohérentes entre elles}

\begin{align*}
  (\alpha-\beta)^{p}\,b
   &= -\struck{\alpha^{p^{2}+p}} + \beta^{p}\alpha^{p^{2}}
      + \struck{\alpha^{p^{2}+p}} - \alpha^{p}\beta^{p^{2}}\\
   &= \alpha^{p}\beta^{p}\bigl(\alpha^{p^{2}-p} - \beta^{p^{2}-p}\bigr)
\end{align*}
\note{les deux termes $\alpha^{p^{2}+p}$ sont barrés d'un même trait oblique,
qui marque leur annulation}

\[
  b = \Bigl(\frac{\alpha\beta}{\alpha-\beta}\Bigr)^{p}
      \bigl(\alpha^{p^{2}-p} - \beta^{p^{2}-p}\bigr)
\]
\note{entre $\alpha$ et $\beta$, au numérateur, un signe raturé qui ne se lit
pas}

\[
  (\alpha-\beta)^{p}X^{p^{2}} \;\ill{}\;
  (\alpha-\beta)^{p^{2}}X^{p}
  + (\alpha\beta)^{p}\bigl(\alpha^{p^{2}-p} - \beta^{p^{2}-p}\bigr)
\]
\note{le signe qui suit $X^{p^{2}}$ est une correction à l'encre lourde, non
lisible ; le dernier terme est écrit sans le facteur $X$}

\[
  (\alpha^{p}-\beta^{p})X^{p^{2}}
  - (\alpha^{p^{2}}-\beta^{p^{2}})X^{p}
  + \alpha^{p}\beta^{p}\bigl(\alpha^{p^{2}-p} - \beta^{p^{2}-p}\bigr)
\]
\note{la page porte une parenthèse ouvrante devant $X^{p^{2}}$, jamais fermée,
non reproduite ; trois accolades soulignent les trois coefficients, qui seront
nommés $r$, $s$, $t$ plus bas}

\[
  \lambda(f^{p}-g^{p})
  \qquad
  -\lambda(f^{p^{2}}-g^{p^{2}})
  \qquad
  \lambda(f^{p^{2}}g^{p} - g^{p^{2}}f^{p})
\]
\note{les trois mêmes coefficients réécrits avec $f$, $g$ et un facteur
$\lambda$}

\[
  \begin{array}{ccc}
    0-\beta^{p} & 0-\beta^{p^{2}} & 0\\
    \alpha^{p}-0 & \alpha^{p^{2}}-\ill{} & 0
  \end{array}
\]
\note{deux lignes de valeurs au trait fin sous les trois coefficients ; la
dernière entrée de la seconde ligne ne se lit pas}

\[
  \alpha^{p} + \beta^{p} = 0
\]

\[
  X \longmapsto \lambda(f^{p}-g^{p})X^{p^{2}}
  - \lambda(f^{p^{2}}-g^{p^{2}})X^{p}
  + \lambda(f^{p^{2}}g^{p} - g^{p^{2}}f^{p})X
\]
\note{les trois coefficients sont accolés et nommés $r$, $s$, $t$}

\begin{tikzcd}[column sep=small, row sep=small, nodes={font=\scriptsize}]
H \arrow[r, hook] & \mathbb{G}_{a} \arrow[r, "X \mapsto rX^{p^{2}} + sX^{p} + t"] \arrow[d] & \mathbb{G}_{a}\\
 & \mathbb{G}_{a}/H &
\end{tikzcd}

\note{le label de la flèche horizontale s'arrête sur $+\,t$, sans le facteur
$X$}

\[
  (f^{p}g - g^{p}f)X^{p^{2}}
  - (f^{p^{2}}g - g^{p^{2}}f)X^{p}
  + (f^{p^{2}}g^{p} - g^{p^{2}}f^{p})X
\]
\note{dernière ligne du feuillet}

\page{22}
\note{le feuillet est écrit des deux bouts : le tiers supérieur à l'endroit,
les deux tiers inférieurs tête-bêche, séparés par un long trait au crayon. On
donne d'abord le bloc à l'endroit, puis le bloc renversé dans son propre ordre
de lecture}

\[
  f \qquad g
\]

\begin{align*}
  f(a) &= 0 & g(a) &= 1\\
  f(b) &= 1 & g(b) &= 0
\end{align*}

donc $f$ et $g$ non cts, lin.\ indép.\ sur \struck{$\mathbb{F}_{p}$}

\uncertain{sps} $f$, $g$ lin.\ indép.\ sur $\mathbb{F}_{p}$ \ill{}

\[
  \mathbb{G}_{a} \supset H_{f}
  \qquad
  \mathbb{G}_{a} \supset H_{g}
\]

\[
  H = H_{f} \times H_{g} \hookrightarrow \mathbb{G}_{a}
\]

\[
  H_{a} \simeq \alpha_{p} \times \mathbb{Z}/p \simeq K \hookrightarrow \mathbb{G}_{a}
\]

\[
  H_{b} \simeq \mathbb{Z}/p \times \alpha_{p} \simeq K \hookrightarrow \mathbb{G}_{a}
\]

\[
  H' \subset \mathbb{G}_{a}
\]

$X \to \struck{X'}$ \ill{}
\note{sous la flèche, trois mots au crayon fin qui ne se lisent pas}

\begin{tikzcd}[column sep=small, row sep=small, nodes={font=\scriptsize}]
H' \subset \mathbb{G}_{a} \arrow[d, "\lambda"]\\
\mathbb{G}_{a}
\end{tikzcd}

$\lambda(a) = \lambda(b) = 0$, $\lambda$ non cte

\[
  H_{a} \simeq H_{b} \;:\; \alpha_{p} \times \mathbb{F}_{p}\alpha \times \mathbb{F}_{p}
\]
\note{le troisième facteur s'arrête sur $\mathbb{F}_{p}$}

\note{ici le trait ; ce qui suit est le bloc renversé, lu depuis le bas du
feuillet}

\[
  r\bigl(\lambda^{p} - f^{p-1}\lambda\bigr)^{p}
  - r^{p}\bigl(\lambda^{p} - f^{p-1}\lambda\bigr)
\]
\note{l'exposant de $r$ au second terme est suivi d'un signe raturé. Dans la
première parenthèse l'exposant de $\lambda$ est mal formé ; il est lu $p$ sur
la seconde parenthèse, qui est nette}

\[
  \Bigl(\frac{\lambda}{f}\Bigr)^{p} - \frac{\lambda}{f}
  \qquad\qquad
  \lambda^{p} - f^{p-1}\lambda
\]
\note{en regard, à droite de la même ligne}

\struck{$\bigl(\lambda^{p} - f^{p-1}\lambda\bigr)\lambda^{p}$}
\note{ligne biffée de deux croix}

\[
  r\lambda^{p^{2}} - f^{p(p-1)}\lambda^{p}
\]

\[
  r = \Bigl(\frac{g}{f}\Bigr)^{p} - f^{p-1}g
\]
\note{le quotient $g/f$ est biffé d'une croix ; corrigé, le second membre est
$g^{p} - f^{p-1}g$}

\[
  \left(
    \frac{\bigl(\frac{\lambda}{f}\bigr)^{p} - \frac{\lambda}{f}}{r}
  \right)^{\uncertain{\nu}}
  = (\quad)
\]
\note{le second membre est une parenthèse laissée vide ; l'exposant est un seul
signe, lu $\nu$}

\[
  a \qquad b
\]

\uncertain{fini plat}
\note{deux mots au crayon fin au-dessus de l'accolade qui joint $a$, $b$ à
l'inclusion suivante}

\[
  H \subset \mathbb{G}_{a},
  \qquad
  H = H_{1} \times H_{2} \times K
\]
\note{la page écrit l'égalité sous forme d'un double trait vertical entre les
deux lignes. Sous $H_{1}$, $H_{2}$ et $K$ descendent trois traits portant deux
lignes d'annotations au trait fin, dont le premier mot semble être une
abréviation en « ét. » suivie d'un chiffre ; elles ne se lisent pas}

\[
  H_{1a} \simeq H_{2b} \simeq \alpha_{p}
\]

\begin{align*}
  f(a) &= 0 & g(b) &= 0\\
  f(b) &= \alpha & g(a) &= \beta
\end{align*}

$\alpha$ et $\beta$ lin.\ indép.\ sur $\mathbb{F}_{p}$

\[
  \alpha_{p} + \mathbb{F}_{p}\alpha + \mathbb{F}_{p}\beta
\]

\note{à partir d'ici le bloc renversé est en deux colonnes séparées d'un trait
vertical : d'abord la colonne de gauche, puis celle de droite}

\[
  H' \subset \mathbb{G}_{a,X'}
\]

\[
  i_{a}(H'_{a}) = i_{b}(H'_{b})
  = \alpha_{p} + \mathbb{F}_{p}\alpha + \mathbb{F}_{p}\beta \subset \mathbb{G}_{a}
\]

\[
  H \subset \mathbb{G}_{a,X}
\]

\[
  i_{a}(H_{a}) = \alpha_{p} + \mathbb{F}_{p}\alpha + \mathbb{F}_{p}\beta
\]

\[
  \mathbb{G}_{a}
\]
\note{$\mathbb{G}_{a}$ est porté seul sur la ligne suivante, sans signe qui le
relie à ce qui précède}

\note{colonne de droite}

$f$, $g$, $h$, lin.\ \ill{} ind.\ sur $\mathbb{F}_{p}$ \ill{}
\note{une accolade horizontale coiffe $f$, $g$ ; à droite, entre accolade,
trois lignes de mots suivies de « $a$, $b$ », qui ne se lisent pas}

\[
  h(a),\ h(b) \in \mathbb{F}_{p}\alpha + \mathbb{F}_{p}\beta
\]

\struck{$g(a)$, $h(a)$}

\struck{\ill{} $H_{2}$}
\note{deux lignes biffées ; de la seconde ne surnage que $H_{2}$. La première
lettre de la première est ambiguë, $f$ ou $g$}

\begin{align*}
  h(a) &\notin \mathbb{F}_{p}\beta & \text{p.\ ex.\ } h(a) &= \alpha\\
  h(b) &\notin \mathbb{F}_{p}\alpha & h(b) &= \beta
\end{align*}

Ker $p_{2}$

\begin{tikzcd}[column sep=small, row sep=small, nodes={font=\scriptsize}]
H \arrow[r, hook] & \mathbb{G}_{a} \arrow[r, "f"] \arrow[dr, "u"'] & \mathbb{G}_{a}\\
 & & \mathbb{G}_{a} \arrow[u, "?"]
\end{tikzcd}

\note{la flèche verticale porte un point d'interrogation, et un trait en
pointillé la longe. À droite du $\mathbb{G}_{a}$ du bas, au crayon, un signe
puis $(H_{2})$ qui ne se lisent pas}

\[
  \mathbb{G}_{a,X'} / H_{1} \times H_{2}
\]

\[
  \lambda \longmapsto
  \Bigl(\frac{\lambda}{f}\Bigr)^{p} - \Bigl(\frac{\lambda}{f}\Bigr)
\]
\note{les deux dernières lignes sont au crayon, plus fines que le reste}

\end{document}
